\documentclass[11pt]{article}
\usepackage{mathblatt}
% gesetzt von werkzeuge/setzer.py (Sorte selbst) aus dem Lernweg FA5
% und den Bankzeilen; Muster bau/proben/2026-10-09/hypotenuse-muster.
\usepackage{enumitem}
\usepackage{adjustbox,varwidth}
\newgeometry{a4paper,margin=18mm,top=15mm,bottom=20mm,footskip=9mm}
\pagestyle{fancy}\fancyhf{}\renewcommand{\headrulewidth}{0pt}
\fancyfoot[L]{\footnotesize\color{mbgrau}FA5 \textperiodcentered{} Lösungen \textperiodcentered{} Zuordnungstypen erkennen und anwenden}
\fancyfoot[R]{\footnotesize\color{mbgrau}\thepage}
\setlength{\parskip}{3pt}
\renewcommand{\kreuz}[1]{\mbox{$\square$\ #1}\quad}
\newcommand{\kreuzl}[1]{\par\hangindent1.4em\hangafter1\noindent$\square$\ #1}
\newcommand{\aufg}[3]{\par\noindent\makebox[0pt][r]{#3}\begin{minipage}[t]{\linewidth}#1\end{minipage}\par}
\newcommand{\aufgg}[4]{\par\noindent\makebox[0pt][r]{#4}\begin{minipage}[t]{0.6\linewidth}\vspace{0pt}#1\end{minipage}\hfill
  \begin{minipage}[t]{0.37\linewidth}\vspace{0pt}\centering #2\end{minipage}\par}
\newcommand{\nr}[1]{\makebox[7mm][l]{\textbf{#1}}}
\newcommand{\lsg}[3]{\par\noindent\begin{minipage}[t]{9mm}\textbf{#1}\end{minipage}%
  \begin{minipage}[t]{52mm}\raggedright\bfseries\boldmath #2\end{minipage}\hspace{3mm}%
  \begin{minipage}[t]{\dimexpr\linewidth-64mm\relax}\raggedright\small #3\end{minipage}\par\vspace{4pt}}
\newlength{\kr}
\makeatletter
\newcommand{\karomerk}[2]{\protected@write\@auxout{}{\string\karoseite{#1}{#2}{\thepage}}}
\newcommand{\karoseite}[3]{\expandafter\xdef\csname karo@#1@#2\endcsname{#3}}
\newcommand{\karohinweis}[3]{\karomerk{h}{#1}\ifcsname karo@k@#1\endcsname
  \ifnum\csname karo@k@#1\endcsname=0\csname karo@h@#1\endcsname\relax#2\else#3\fi
  \else#3\fi}
\makeatother
\begin{document}

\noindent{\Large\bfseries Lösungen}\hspace{1em}{\color{mbgrau}Zuordnungstypen erkennen und anwenden}\par\smallskip
\noindent{\small Links das Ergebnis, rechts der Weg in kurzen Schritten. Stimmt dein Ergebnis nicht, such den ersten Schritt, der anders ist.}\par
\par\Needspace{25mm}\vspace{6pt}\noindent\colorbox{black!12}{\makebox[7mm]{\bfseries\strut A}}\hspace{2mm}{\bfseries Tabelle und Graph prüfen}\par\vspace{4pt}
\lsg{1.}{proportional}{$6 : 2 = 12 : 4 = 15 : 5 = 24 : 8 = 3$}
\lsg{2.}{antiproportional}{Quotienten $30$, $7{,}5$, … nicht gleich; Produkte $1 \cdot 30 = 2 \cdot 15 = 3 \cdot 10 = 5 \cdot 6 = 30$}
\lsg{3.}{keins von beiden}{Quotienten $5 : 1 = 5$, $7 : 2 = 3{,}5$; Produkte $1 \cdot 5 = 5$, $2 \cdot 7 = 14$; beide nicht gleich}
\lsg{4.}{(a) keins \ (b) proportional \ (c) antiproportional}{(a) Gerade, aber nicht durch den Ursprung; (b) Gerade durch den Ursprung; (c) fallende Kurve}
\lsg{5.}{keins von beiden; kein Dreisatz}{Keins von beiden; Quotienten $9 : 2 = 4{,}5$, $13 : 4 = 3{,}25$ – nicht gleich; Produkte $2 \cdot 9 = 18$, $4 \cdot 13 = 52$ – nicht gleich. Also kein Dreisatz (der Verleih nimmt eine Grundgebühr von 5\,€ und 2\,€ je Stunde).}
\par\Needspace{25mm}\vspace{6pt}\noindent\colorbox{black!12}{\makebox[7mm]{\bfseries\strut B}}\hspace{2mm}{\bfseries Am Text erkennen, dann rechnen}\par\vspace{4pt}
\lsg{1.}{3{,}20\,€}{proportional: $2 : 5 = 0{,}40$; $0{,}40 \cdot 8 = 3{,}20$\,€}
\lsg{2.}{6 Stunden}{antiproportional: $3 \cdot 8 = 24$; $24 : 4 = 6$\,h}
\lsg{3.}{16 Minuten}{antiproportional (doppelt so schnell, halb so lange): $12 \cdot 20 = 240$; $240 : 15 = 16$\,min}
\lsg{4.}{33 Minuten}{proportional (doppelt so weit, doppelt so lange): $22 : 4 = 5{,}5$; $5{,}5 \cdot 6 = 33$\,min}
\lsg{5.}{Nein, ihr Anteil ist 105\,€; es fehlen 5\,€.}{Nein; antiproportional: $6 \cdot 140 = 840$; $840 : 8 = 105$\,€; $105 - 100 = 5$\,€ fehlen.}
\par\Needspace{25mm}\vspace{6pt}\noindent\colorbox{black!12}{\makebox[7mm]{\bfseries\strut C}}\hspace{2mm}{\bfseries Mit der Rate rechnen: Kosten, Tempo, Dauer}\par\vspace{4pt}
\lsg{1.}{4{,}20\,€}{$12 \cdot 35 = 420$\,ct $= 4{,}20$\,€}
\lsg{2.}{6 Minuten}{$900 : 2{,}5 = 360$\,s; $360 : 60 = 6$\,min}
\lsg{3.}{16\,km/h}{$15$\,min $\cdot 4 = 60$\,min; $4$\,km $\cdot 4 = 16$\,km in 1 Stunde}
\lsg{4.}{45 Minuten}{$18 : 24 = 0{,}75$\,h; $0{,}75 \cdot 60 = 45$\,min}
\lsg{5.}{Das stimmt: 324\,€ gespart.}{$12\,000 : 100 = 120$; $120 \cdot 6 = 720$\,l; $120 \cdot 4{,}5 = 540$\,l; $720 - 540 = 180$\,l; $180 \cdot 1{,}80 = 324$\,€ $> 300$\,€}
\par\Needspace{25mm}\vspace{6pt}\noindent\colorbox{black!12}{\makebox[7mm]{\bfseries\strut T}}\hspace{2mm}{\bfseries Probetest}\par\vspace{4pt}
\lsg{1.}{keins von beiden}{Quotienten $0{,}7$, $0{,}45$; Produkte $70$, $180$; beide nicht gleich}
\lsg{2.}{8 Minuten}{proportional: $45 : 3 = 15$ Seiten je Minute; $120 : 15 = 8$\,min}
\lsg{3.}{10 Tage}{antiproportional: $8 \cdot 15 = 120$; $120 : 12 = 10$ Tage}
\lsg{4.}{2\,h 15\,min}{$5400 : 40 = 135$\,min; $135 = 2 \cdot 60 + 15$, also 2\,h 15\,min}
\lsg{5.}{25 Personen: Bahn (Bus 24\,€); 30 Personen: Bus (20\,€)}{Bei 25: Bahn; Bus $600 : 25 = 24$\,€ $> 22$\,€. Bei 30: Bus; $600 : 30 = 20$\,€ $< 22$\,€.}
\end{document}
